Using the calibrated bounds in Eq.~\eqref{eq:calibrated_energy_bound}, define the minimum calibrated cost bound for each trajectory type as
\begin{equation}
    \overline c_i^\tau:=\min_{k\in\{1,\ldots,K\}}\overline c_{i,k}^\tau,
    \qquad
    k_i^\tau\in\arg\min_k \overline c_{i,k}^\tau,
    \label{eq:minimum_calibrated_cost}
\end{equation}
where station ties are resolved by a fixed rule.
The station-reachability certificate and energy margin are then defined~as
\begin{equation}
    h_{e,i}=V_i-V_{\min}-\overline c_i^{\mathrm{sta}},
    \qquad
    \mu_i=V_i-V_{\min}-\overline c_i^{\mathrm{task}}.
    \label{eq:calibrated_energy_quantities}
\end{equation}
If $h_{e,i}\geq0$ and the station-cost bound is not violated, then
\begin{equation}
    V_i - V_{\min}-c_{i, k_i^\mathrm{sta}}^{\mathrm{sta}}
    \geq V_i - V_{\min} -\overline{c}_{i, k_i^{\mathrm{sta}}}^{\mathrm{sta}} = h_{e, i} \geq 0.
    \label{eq:conditional_station_reachability}
\end{equation}
Thus, $h_{e,i}\geq0$ provides a conditional certificate that robot $i$ can reach the selected station without falling below $V_{\min}$.
At runtime, $h_{e,i}$ is evaluated as an online station-reachability certificate, but is not imposed as a control constraint.

For right-of-way coordination, a larger $\mu_i$ indicates greater flexibility to accommodate the additional cost of yielding.
For an interacting pair $(i,j)$, define
\begin{equation}
    i\rightarrow j
    \quad\Longleftrightarrow\quad
    \mu_i>\mu_j
    \;\;\text{or}\;\;
    \bigl(\mu_i=\mu_j \text{ and } i>j\bigr),
    \label{eq:energy_precedence}
\end{equation}
where $i\rightarrow j$ means that robot $i$ yields to robot $j$.
Because the margins and robot indices define a strict total order, each interacting pair receives exactly one yielding direction.
Any subset of pairwise relations induced by this order is therefore acyclic.
The robot index is used only when the calibrated margins are equal; all non-tied decisions are determined by the energy-sufficiency margins.
The ordering in Eq.~\eqref{eq:energy_precedence} satisfies the energy-guided precedence requirement of Problem~\ref{prob:main}.
